\documentclass[11pt,a4paper]{article}

% ──────────────────────────────────────────────
% Preamble
% ──────────────────────────────────────────────
\usepackage[utf8]{inputenc}
\usepackage[T1]{fontenc}
\usepackage{lmodern}
\usepackage[margin=1in]{geometry}
\usepackage{amsmath,amssymb,amsthm}
\usepackage{booktabs}
\usepackage{graphicx}
\usepackage{xcolor}
\usepackage{hyperref}
\usepackage{float}
\usepackage{enumitem}
\usepackage[most]{tcolorbox}
\usepackage{fancyhdr}
\usepackage{titlesec}
\usepackage{microtype}
\usepackage{datetime2}

% ── Colors ───────────────────────────────────
\definecolor{uscCardinal}{RGB}{153,0,0}
\definecolor{uscGold}{RGB}{255,204,0}
\definecolor{darkblue}{RGB}{0,51,102}
\definecolor{insightbg}{RGB}{255,248,220}
\definecolor{insightframe}{RGB}{204,153,0}
\definecolor{bugred}{RGB}{220,50,50}
\definecolor{fixgreen}{RGB}{40,160,60}
\definecolor{codebg}{RGB}{245,245,245}

% ── tcolorbox environments ───────────────────
\newtcolorbox{keyinsight}[1][]{%
  enhanced, breakable,
  colback=insightbg, colframe=insightframe,
  fonttitle=\bfseries\large,
  title={Key Insight},
  #1
}

\newtcolorbox{bugbox}[1][]{%
  enhanced, breakable,
  colback=red!3, colframe=bugred!70,
  fonttitle=\bfseries,
  title={#1}
}

\newtcolorbox{fixbox}[1][]{%
  enhanced, breakable,
  colback=green!3, colframe=fixgreen!70,
  fonttitle=\bfseries,
  title={#1}
}

\newtcolorbox{configbox}[1][]{%
  enhanced, breakable,
  colback=codebg, colframe=gray!50,
  fonttitle=\bfseries,
  title={#1}
}

% ── Headers ──────────────────────────────────
\pagestyle{fancy}
\fancyhf{}
\fancyhead[L]{\small\textcolor{gray}{Offlineness in GRPO --- Experiment Log}}
\fancyhead[R]{\small\textcolor{gray}{\thepage}}
\renewcommand{\headrulewidth}{0.4pt}

% ── Section formatting ───────────────────────
\titleformat{\section}
  {\Large\bfseries\color{darkblue}}{\thesection}{1em}{}[\vspace{-0.5em}\textcolor{darkblue}{\rule{\textwidth}{0.5pt}}]
\titleformat{\subsection}
  {\large\bfseries\color{darkblue!80}}{\thesubsection}{1em}{}

% ── Hyperref setup ───────────────────────────
\hypersetup{
  colorlinks=true,
  linkcolor=darkblue,
  citecolor=darkblue,
  urlcolor=uscCardinal
}

% ── Document info ────────────────────────────
\title{%
  \vspace{-1cm}
  {\LARGE\bfseries\color{darkblue} Offlineness in GRPO Training\\[0.3em] Without KL Penalty}\\[0.8em]
  {\Large Experiment Log --- SLIME Framework}\\[0.5em]
  {\normalsize\color{gray} Living Document --- Last Updated: \today}
}
\author{%
  Yanxin Peng\\
  \small University of Southern California\\
  \small \texttt{yanxinpe@usc.edu}
}
\date{}

% ══════════════════════════════════════════════
\begin{document}
\maketitle
\thispagestyle{fancy}

\tableofcontents
\newpage

% ══════════════════════════════════════════════
\section{Research Motivation \& Background}
\label{sec:motivation}
% ══════════════════════════════════════════════

Group Relative Policy Optimization (GRPO) has emerged as a leading algorithm for
post-training large language models on mathematical reasoning tasks. Unlike
standard PPO which constrains the policy via a KL penalty term against a
reference policy, GRPO in its purest form sets the KL coefficient to zero
($\lambda_{\mathrm{KL}} = 0$), relying solely on importance-weight clipping to
prevent excessive policy updates.

This design choice raises a fundamental question: \textbf{how off-policy does
the training data become over the course of training, and does this
off-policyness threaten training stability?}

In on-policy reinforcement learning, the data used to compute policy gradients
is generated by the current policy $\pi_\theta$. In practice, however, there is
always a gap between the data-generating policy $\pi_{\mathrm{old}}$ (the policy
at rollout time) and the policy $\pi_\theta$ at training time. This gap grows
with each gradient step taken on the same batch of rollout data. We term this
gap \emph{offlineness}.

Off-policy data degrades the importance sampling estimates that underpin policy
gradient methods. When the importance weights $w_i = \pi_\theta(a_i \mid s_i) /
\pi_{\mathrm{old}}(a_i \mid s_i)$ deviate significantly from 1, the effective
sample size drops, gradient variance increases, and the policy may collapse.

\subsection{Experimental Setup}

\begin{configbox}{Training Configuration}
\begin{tabular}{@{}ll@{}}
  \textbf{Model} & Qwen2.5-Math-1.5B \\
  \textbf{Training data} & DAPO-math-17k (17{,}000 math prompts) \\
  \textbf{Framework} & SLIME v0.2.1 (THUDM/Tsinghua) \\
  \textbf{Backend} & Megatron-LM (tensor parallel) + SGLang (rollout) \\
  \textbf{GPUs} & 4$\times$ RTX A6000 (2 training + 2 rollout) \\
  \textbf{RL algorithm} & GRPO with $\lambda_{\mathrm{KL}} = 0$ \\
  \textbf{PPO epochs} & 1 (single pass over each rollout batch) \\
  \textbf{Clipping} & $\epsilon_{\mathrm{clip}} = 0.2$, $\epsilon_{\mathrm{clip,high}} = 0.28$ \\
  \textbf{Batch size} & Global 128, rollout batch 16, 8 samples/prompt \\
  \textbf{Learning rate} & $10^{-6}$ (constant schedule) \\
  \textbf{Infrastructure} & USC CARC Endeavour cluster, NLP partition \\
\end{tabular}
\end{configbox}


% ══════════════════════════════════════════════
\section{Offlineness Metrics}
\label{sec:metrics}
% ══════════════════════════════════════════════

We instrument SLIME's training loop with four metrics that quantify offlineness
from complementary perspectives. All metrics are computed per microbatch during
the policy loss computation, using per-token importance weights:
\[
  w_i = \frac{\pi_\theta(a_i \mid s_i)}{\pi_{\mathrm{old}}(a_i \mid s_i)}
      = \exp\!\bigl(\log\pi_\theta(a_i \mid s_i) - \log\pi_{\mathrm{old}}(a_i \mid s_i)\bigr)
\]
where $\pi_{\mathrm{old}}$ denotes the policy at \emph{rollout time} (data
generation), not the start of the current PPO epoch.

\subsection{Effective Sample Size Ratio (ESS)}

The ESS measures how efficiently the importance-weighted samples approximate
on-policy data:
\begin{equation}
  \mathrm{ESS} = \frac{\bigl(\sum_{i=1}^{N} w_i\bigr)^2}{N \cdot \sum_{i=1}^{N} w_i^2}
  \label{eq:ess}
\end{equation}
The ratio is normalized to $[0, 1]$: a value of 1.0 means perfectly on-policy
(all weights equal), while values approaching 0 indicate severe off-policy
degradation where a few samples dominate the gradient estimate.

\subsection{Mismatch KL ($\chi^2$ Divergence Proxy)}

We use a second-order approximation to the KL divergence:
\begin{equation}
  D_{\mathrm{mismatch}} = \frac{1}{2}\,\mathbb{E}\bigl[(w_i - 1)^2\bigr]
  \label{eq:mismatch_kl}
\end{equation}
This is equivalent to the $k_3$ KL estimator and penalizes large importance
weights quadratically. It is more sensitive to outlier weights than the standard
KL divergence.

\subsection{Clip Fraction}

The fraction of tokens whose importance weight falls outside a reference
interval:
\begin{equation}
  \mathrm{clip\_frac} = \frac{1}{N}\sum_{i=1}^{N}
    \mathbf{1}\!\bigl[w_i \notin [0.8,\; 1.2]\bigr]
  \label{eq:clip_frac}
\end{equation}
This directly measures the proportion of tokens that have ``drifted'' beyond an
acceptable range.

\subsection{Maximum Ratio}

\begin{equation}
  w_{\max} = \max_{i} \; w_i
  \label{eq:max_ratio}
\end{equation}

The worst-case importance weight across all tokens in the microbatch. A single
extreme weight can destabilize the gradient estimate even if the average weight
is close to 1.


% ══════════════════════════════════════════════
\section{Implementation}
\label{sec:implementation}
% ══════════════════════════════════════════════

We implement offlineness monitoring directly in SLIME's policy loss function.
The implementation consists of three components:

\begin{enumerate}[leftmargin=*]
  \item \textbf{Metric computation} (\texttt{compute\_offlineness\_metrics()} in
        \texttt{loss.py}, lines 391--427): Takes current and reference
        log-probabilities, applies loss masks to select valid tokens, computes
        all four metrics from the importance weight tensor.

  \item \textbf{Integration with training loop} (\texttt{policy\_loss\_function()}
        in \texttt{loss.py}, lines 697--741): Metrics are computed at the end of
        each microbatch's forward pass and inserted into the loss dictionary for
        logging.

  \item \textbf{CLI flag} (\texttt{arguments.py}): The
        \texttt{-{}-enable-offlineness-metrics} flag gates all computation,
        ensuring zero overhead when monitoring is disabled.
\end{enumerate}

Critically, the reference log-probabilities must come from
\texttt{batch["rollout\_log\_probs"]} (the policy at data generation time), not
\texttt{batch["log\_probs"]} (the policy at the start of the current PPO epoch).
The latter measures only within-epoch drift and is near-zero by construction.


% ══════════════════════════════════════════════
\section{Bug Discovery \& Fixes}
\label{sec:bugs}
% ══════════════════════════════════════════════

Bringing up the offlineness metrics revealed two bugs in the instrumentation
code. Both were discovered, diagnosed, and fixed within a single day
(February~26, 2026). We document them in detail as they illustrate subtle
pitfalls in Megatron's metric aggregation pipeline.

% ── Bug 1 ────────────────────────────────────
\subsection{Bug 1: Wrong Reference Log-Probabilities}

\begin{bugbox}{Bug 1 --- Wrong Reference Log-Probs (Discovered Feb 26)}
\textbf{Symptom}: All offlineness metrics showed near-zero values
  ($D_{\mathrm{mismatch}} = 0$, $\mathrm{clip\_frac} = 0$) in every training step.\\[0.3em]
\textbf{Root cause}: The code compared the current policy's log-probs against
  \texttt{batch["log\_probs"]}, which stores the log-probabilities from the
  \emph{start of the current PPO epoch}---nearly identical to the current policy
  within one epoch. This measured within-epoch drift (effectively zero) rather
  than true offlineness from the rollout policy.\\[0.3em]
\textbf{Impact}: The metrics appeared to show a perfectly on-policy regime,
  masking any actual drift.
\end{bugbox}

\begin{fixbox}{Fix (Commit \texttt{80a3cc6})}
Changed the reference from \texttt{batch["log\_probs"]} to
  \texttt{batch["rollout\_log\_probs"]}, which stores the per-token
  log-probabilities captured at data generation time. After this fix,
  $D_{\mathrm{mismatch}}$ and $\mathrm{clip\_frac}$ became non-zero, but
  $\mathrm{ESS}$ still displayed anomalous values ($\approx 0.19$).
\end{fixbox}

% ── Bug 2 ────────────────────────────────────
\subsection{Bug 2: Metric Reduction Error}

\begin{bugbox}{Bug 2 --- Megatron Reduction Pipeline (Discovered Feb 26)}
\textbf{Symptom}: After fixing Bug~1, the WandB dashboard displayed
  $\mathrm{ESS} \approx 0.19$. However, debug prints inserted into
  \texttt{compute\_offlineness\_metrics()} showed per-microbatch values of
  $\mathrm{ESS} \approx 0.999$. This mathematical contradiction prompted a
  deeper investigation.\\[0.3em]
\textbf{Root cause}: Megatron's training pipeline aggregates metrics by:
  \begin{enumerate}[leftmargin=1.5em, topsep=0.2em]
    \item Packing all metric values into a flat tensor per microbatch:
          $[\texttt{num\_samples},\; m_1,\; m_2,\; \ldots]$
    \item \emph{Summing} across microbatches
    \item Dividing each metric by the total \texttt{num\_samples}
  \end{enumerate}
  This pipeline correctly handles sum-reducible quantities (like total loss or
  sum-of-sample-means). However, our offlineness metrics are \emph{scalar
  averages} (e.g., $\mathrm{ESS} = 0.999$). Inserting a raw scalar into the
  pipeline yields:
  \[
    \frac{0.999 \times 16\;\text{microbatches}}{128\;\text{total samples}}
    \;=\; 0.125
  \]
  which is close to the observed $\approx 0.19$.\\[0.3em]
\textbf{Impact}: All four offlineness metrics were reported at $\sim\!1/8$ of
  their true values, making the policy appear dangerously off-policy when it was
  in fact nearly on-policy.
\end{bugbox}

\begin{fixbox}{Fix (Commit \texttt{ac1de44})}
Before inserting offlineness metrics into \texttt{reported\_loss}, we
  pre-multiply each metric by the microbatch's \texttt{num\_samples}:
  \[
    \texttt{reported\_loss[key]} = \texttt{value} \times \texttt{num\_samples}
  \]
  After the pipeline sums across microbatches and divides by
  \texttt{total\_samples}, the correct weighted average is recovered. Verified:
  ESS now reports $0.996$--$0.999$ consistently.
\end{fixbox}


% ══════════════════════════════════════════════
\section{Baseline Results}
\label{sec:baseline}
% ══════════════════════════════════════════════

With both bugs fixed (commit \texttt{ac1de44}), we ran the baseline
configuration for 34 training steps (steps 600--633, resuming from a checkpoint).
Table~\ref{tab:baseline} summarizes the corrected offlineness metrics alongside
standard training indicators.

\begin{table}[H]
\centering
\caption{Baseline metrics over steps 600--633 (34 steps). All offlineness
metrics use the corrected reduction pipeline.}
\label{tab:baseline}
\begin{tabular}{@{}lcccc@{}}
\toprule
\textbf{Metric} & \textbf{Mean} & \textbf{Std} & \textbf{Min} & \textbf{Max} \\
\midrule
\multicolumn{5}{@{}l}{\textit{Offlineness Metrics}} \\
\quad ESS Ratio            & 0.9988 & 0.0007 & 0.9965 & 0.9992 \\
\quad Mismatch KL          & 0.0006 & 0.0004 & 0.0004 & 0.0019 \\
\quad Clip Fraction (0.2)  & 0.0028 & 0.0003 & 0.0023 & 0.0034 \\
\quad Max Ratio            & 1.72   & 0.24   & 1.43   & 2.30   \\
\midrule
\multicolumn{5}{@{}l}{\textit{Training Metrics}} \\
\quad Policy Loss          & $\approx 0$ & --- & --- & --- \\
\quad PG Clip Fraction     & 0.0    & 0.0    & 0.0    & 0.0    \\
\quad Entropy Loss         & 0.22   & 0.02   & 0.19   & 0.28   \\
\quad Gradient Norm        & 0.11   & 0.02   & 0.08   & 0.14   \\
\midrule
\multicolumn{5}{@{}l}{\textit{Rollout Statistics}} \\
\quad Raw Reward           & 0.19   & 0.05   & 0.13   & 0.30   \\
\quad Response Length      & 984    & 56     & 828    & 1058   \\
\quad Truncation Rate      & 0.08   & 0.02   & 0.05   & 0.13   \\
\bottomrule
\end{tabular}
\end{table}

Several observations emerge from the baseline:

\begin{itemize}[leftmargin=*]
  \item \textbf{Near-perfect on-policyness.} The ESS ratio of $0.9988$ indicates
    that virtually all importance weights are close to~1. The importance sampling
    estimator is as efficient as an on-policy estimator.

  \item \textbf{Negligible mismatch.} The mismatch KL of $6 \times 10^{-4}$ and
    clip fraction of $0.28\%$ confirm minimal divergence between the rollout and
    training policies.

  \item \textbf{Per-token extremes exist but are rare.} The maximum ratio reaches
    up to 2.3, meaning some individual tokens see their probability double. But
    these extremes are isolated ($<\!0.3\%$ of tokens) and do not significantly
    affect the aggregate ESS.

  \item \textbf{Zero effective policy gradient.} The policy loss is numerically
    zero ($\approx 10^{-9}$), and the PPO clip fraction is exactly 0. With
    \texttt{kl\_coef=0}, the PPO clip fraction compares \emph{within-epoch}
    ratios (which are trivially near~1 with a single PPO epoch), not rollout
    ratios.

  \item \textbf{Moderate reward level.} The raw reward of 19\% (solve rate)
    suggests the model is in an active learning phase, not fully converged.
\end{itemize}


% ══════════════════════════════════════════════
\section{Key Insight \& Revised Research Direction}
\label{sec:insight}
% ══════════════════════════════════════════════

\begin{keyinsight}
The original hypothesis was that removing the KL penalty ($\lambda_{\mathrm{KL}}
= 0$) would lead to dangerous off-policy drift, as measured by ESS collapse.

\textbf{Finding:} With the standard GRPO configuration (1~PPO epoch per rollout,
$\mathrm{lr} = 10^{-6}$, $\epsilon_{\mathrm{clip}} = 0.2$), the policy remains
\emph{almost perfectly on-policy} ($\mathrm{ESS} \approx 0.999$). The single
PPO epoch and conservative learning rate ensure that the policy barely moves
between data generation and training.

This means the current setup is inherently safe from off-policy degradation---but
also potentially \textbf{underutilizing the rollout data}. Each expensive rollout
batch is used for only one gradient step before being discarded.
\end{keyinsight}

This finding reframes the research question. Rather than asking ``does
offlineness harm GRPO training?'' (answer: not in the default configuration), we
now ask:

\begin{enumerate}[leftmargin=*]
  \item \textbf{Can we deliberately induce offlineness} by increasing data reuse
    (more PPO epochs), raising the learning rate, or reducing rollout frequency?

  \item \textbf{What is the ESS threshold} below which training becomes unstable
    (reward collapse, entropy collapse, gradient explosion)?

  \item \textbf{Can we build an adaptive mechanism} that monitors ESS in real
    time and adjusts PPO epochs or learning rate to stay in a ``sweet spot''---
    maximizing data efficiency while maintaining stability?
\end{enumerate}


% ══════════════════════════════════════════════
\section{Proposed Ablation Experiments}
\label{sec:ablations}
% ══════════════════════════════════════════════

To answer the revised research questions, we propose the following ablation
study. All experiments share the baseline configuration except for the varied
parameter.

\subsection{Experiment 1: PPO Epoch Sweep}

\textbf{Rationale.} Each additional PPO epoch reuses the same rollout data for
another gradient step, causing the training policy to drift further from the
rollout policy. This is the most natural way to control offlineness.

\begin{table}[H]
\centering
\caption{PPO epoch ablation design.}
\label{tab:ppo_sweep}
\begin{tabular}{@{}lcl@{}}
\toprule
\textbf{Variant} & \textbf{PPO Epochs} & \textbf{Expected Effect} \\
\midrule
Baseline     & 1  & ESS $\approx 0.999$ (near on-policy) \\
Moderate     & 3  & ESS slightly decreased \\
Aggressive   & 5  & ESS noticeably decreased \\
Extreme      & 10 & ESS potentially in danger zone \\
\bottomrule
\end{tabular}
\end{table}

\subsection{Experiment 2: Learning Rate Sweep}

\textbf{Rationale.} A higher learning rate causes larger policy updates per
gradient step, amplifying the per-step drift.

\begin{table}[H]
\centering
\caption{Learning rate ablation design.}
\label{tab:lr_sweep}
\begin{tabular}{@{}lcl@{}}
\toprule
\textbf{Variant} & \textbf{Learning Rate} & \textbf{Expected Effect} \\
\midrule
Baseline     & $10^{-6}$ & ESS $\approx 0.999$ \\
$10\times$   & $10^{-5}$ & Faster drift per step \\
$100\times$  & $10^{-4}$ & Potentially unstable \\
\bottomrule
\end{tabular}
\end{table}

\subsection{Experiment 3: Combined (PPO Epochs + Learning Rate)}

\textbf{Rationale.} The interaction between PPO epochs and learning rate may be
nonlinear. More epochs with a higher learning rate could push the system into
a qualitatively different regime.


% ══════════════════════════════════════════════
\section{Ablation Experiment 1: PPO Epoch Sweep}
\label{sec:ppo_epoch_results}
% ══════════════════════════════════════════════

\subsection{Implementation}

SLIME's original codebase does not support true PPO epochs---it makes a single
pass over each rollout batch. We implemented multi-epoch PPO training by adding a
loop around the \texttt{train()} call in \texttt{actor.py:train\_actor()}
(commit \texttt{6205e02}):

\begin{enumerate}[leftmargin=*]
  \item A new CLI argument \texttt{-{}-ppo-epochs} (default~1) controls the
    number of training passes per rollout.
  \item For each epoch after the first, the data iterators are reset and
    \texttt{log\_probs} are \emph{recomputed} under the updated policy.
    This ensures the PPO importance ratio $w_i = \pi_\theta / \pi_{\mathrm{old}}$
    compares against the epoch-start policy (not the first epoch's policy),
    keeping ratios bounded.
  \item Advantages are \emph{not} recomputed between epochs. For GRPO with
    $\lambda_{\mathrm{KL}} = 0$, advantages depend only on rewards (which are
    fixed for the rollout), so recomputation is unnecessary.
\end{enumerate}

Note the distinction between ``PPO epochs'' and SLIME's existing
\texttt{-{}-num-steps-per-rollout} parameter. The latter \emph{splits} rollout
data into non-overlapping chunks (e.g., 128 samples into 4 batches of 32),
providing only one gradient update per sample. True PPO epochs loop over the
\emph{same} full data $K$ times, providing $K$ gradient updates per sample.

\subsection{Experiment Configuration}

All PPO epoch experiments share the baseline configuration (Section~\ref{sec:baseline})
except for the \texttt{-{}-ppo-epochs} parameter and WandB group name. Each
experiment resumes from the step~599 checkpoint.

\begin{table}[H]
\centering
\caption{PPO epoch sweep experiment configuration.}
\label{tab:ppo_sweep_config}
\begin{tabular}{@{}lccl@{}}
\toprule
\textbf{Run} & \textbf{PPO Epochs} & \textbf{WandB Group} & \textbf{Status} \\
\midrule
Baseline     & 1  & \texttt{ppo-epochs-1-baseline} & Complete (steps 600--633) \\
Moderate     & 3  & \texttt{ppo-epochs-3}          & Running (SLURM 2769138)   \\
Aggressive   & 5  & \texttt{ppo-epochs-5}          & Running (SLURM 2769153)   \\
Extreme      & 10 & \texttt{ppo-epochs-10}         & Planned                   \\
\bottomrule
\end{tabular}
\end{table}

\subsection{Early Results: PPO Epochs = 3}

The PPO epoch=3 experiment began running on Feb~27, 2026 (node~c05-01).
Initial observations from steps 600--603:

\begin{table}[H]
\centering
\caption{PPO epoch=3 metrics (first epoch of each rollout), steps 600--603.}
\label{tab:ppo3_early}
\begin{tabular}{@{}lcccc@{}}
\toprule
\textbf{Metric} & \textbf{Step 600} & \textbf{Step 601} & \textbf{Step 602} & \textbf{Step 603} \\
\midrule
ESS Ratio      & 0.9992 & 0.9988 & 0.9987 & 0.9991 \\
Mismatch KL    & 0.00039 & 0.00059 & 0.00068 & 0.00045 \\
Clip Frac (0.2) & 0.0023 & 0.0025 & 0.0025 & 0.0034 \\
Max Ratio       & 1.529  & 1.883  & 1.805  & 1.570  \\
Grad Norm       & 0.110  & 0.107  & 0.107  & 0.118  \\
\bottomrule
\end{tabular}
\end{table}

\begin{keyinsight}
\textbf{Preliminary finding:} All three PPO epochs within the same rollout
produce \emph{nearly identical} offlineness metrics. The ESS ratio remains
$>0.998$ and the grad norm stays at $\sim\!0.11$ across all epochs. This
suggests that at the current learning rate ($10^{-6}$), the per-epoch gradient
updates are too small to cause meaningful policy drift within a single rollout.

The cumulative effect over hundreds of rollouts remains to be observed---with
3$\times$ more gradient updates per rollout, the model may diverge from the
baseline trajectory over time even if per-rollout drift is negligible.
\end{keyinsight}


% ══════════════════════════════════════════════
\section{Experiment Timeline}
\label{sec:timeline}
% ══════════════════════════════════════════════

\begin{table}[H]
\centering
\caption{Chronological log of all experiments.}
\label{tab:timeline}
\begin{tabular}{@{}llp{7cm}@{}}
\toprule
\textbf{Date} & \textbf{Experiment} & \textbf{Outcome} \\
\midrule
Feb 26, 2026 & Buggy baseline (pre-fix) & ESS $\approx 0.19$ (artifact of Bug~2); mismatch\_kl $= 0$ (Bug~1) \\
Feb 26, 2026 & Debug run (Bug~1 fixed) & ESS still $\approx 0.19$; per-microbatch debug prints reveal ESS $\approx 0.999$; Bug~2 identified \\
Feb 26, 2026 & Fixed baseline (both bugs fixed) & ESS $= 0.9988 \pm 0.0007$ over steps 600--633; all metrics validated \\
Feb 27, 2026 & PPO epochs code (commit \texttt{6205e02}) & Implemented \texttt{-{}-ppo-epochs} with log\_probs recomputation \\
Feb 27, 2026 & PPO epoch=3 (SLURM 2769138) & Running on c05-01; early results show ESS $\approx 0.999$ across all 3 epochs \\
Feb 27, 2026 & PPO epoch=5 (SLURM 2769153) & Running on c06-01; awaiting initial results \\
\midrule
\multicolumn{3}{@{}l}{\textit{Upcoming}} \\
TBD & PPO epoch=10 & Planned \\
TBD & Learning rate sweep & --- \\
TBD & Combined ablation (epochs + LR) & --- \\
\bottomrule
\end{tabular}
\end{table}


% ══════════════════════════════════════════════
\section{Next Steps}
\label{sec:next}
% ══════════════════════════════════════════════

\begin{enumerate}[leftmargin=*]
  \item \textbf{Collect PPO epoch sweep data.} Let the epoch=3 and epoch=5
    experiments run for at least 100 rollouts ($\sim$2 hours each) to observe
    cumulative effects on ESS, reward, and entropy.

  \item \textbf{Launch PPO epoch=10 experiment.} The most aggressive data reuse
    setting; expected to show measurable ESS degradation if any.

  \item \textbf{Learning rate sweep.} If PPO epochs alone do not induce
    offlineness, increase the learning rate ($10^{-5}$, $10^{-4}$) to amplify
    per-step drift.

  \item \textbf{Identify the danger zone.} Determine the ESS threshold below
    which training exhibits instability (reward plateau, entropy collapse,
    gradient explosion).

  \item \textbf{Characterize the efficiency--stability tradeoff.} Map reward
    gain vs.\ ESS degradation as a function of PPO epochs and learning rate.

  \item \textbf{Propose adaptive ESS-based early stopping.} If ESS drops below
    a threshold mid-training, stop the PPO inner loop early and trigger a new
    rollout.
\end{enumerate}


% ══════════════════════════════════════════════
% Appendix: Technical Notes
% ══════════════════════════════════════════════
\appendix
\section{Technical Notes}
\label{app:technical}

\subsection{Megatron Metric Reduction Pipeline}

The Megatron training loop processes data in microbatches. For each microbatch,
the loss function returns a dictionary \texttt{reported\_loss} containing metric
tensors. These are packed into a flat tensor:
\[
  \texttt{values} = [\texttt{num\_samples},\; m_1,\; m_2,\; \ldots]
\]
Across microbatches, values are \emph{summed}. An \texttt{all\_reduce} aggregates
across data-parallel ranks. Finally, each metric is computed as:
\[
  \hat{m}_k = \frac{\texttt{values}[k]}{\texttt{values}[0]} \times \texttt{cp\_size}
\]
where $\texttt{values}[0]$ is the total \texttt{num\_samples}. This design
assumes that each $m_k$ was originally a ``sum-of-sample-means'' (as produced by
\texttt{sum\_of\_sample\_mean()} for losses). For scalar-average metrics like
ESS, we must pre-multiply by $\texttt{num\_samples}$ per microbatch to ensure
correctness.

\subsection{Cluster Infrastructure}

Experiments run on the USC CARC Endeavour cluster (\texttt{nlp} partition).
Notable operational details:

\begin{itemize}[leftmargin=*]
  \item Nodes \texttt{c05-05} and \texttt{c06-05} have broken CUDA drivers
    (Error 802); always excluded via \texttt{-{}-exclude}.
  \item The NLP partition supports preemption---jobs can be killed by
    higher-priority users. Experiments must be designed to resume from
    checkpoints.
  \item Job management is fully automated via \texttt{sbatch} scripts that
    include \texttt{pip install -e .} for code synchronization.
\end{itemize}


\end{document}
